\paragraph{Notation}

We use $\lfloor \cdot \rfloor$, $\lceil \cdot \rceil$, and
$\lfloor \cdot \rceil$ to denote rounding down, up, and to the
nearest integer, respectively. When these operations are applied
to a polynomial, we mean performing the corresponding operation
to each coefficient separately. The norm $\|\cdot\|$ denotes the
infinity norm and $\|\cdot\|_{\mathrm{can}}$ denotes the canonical
norm. We denote the reduction of an integer modulo
$t$ by $[\cdot]_t$. This operation can also be applied to
polynomials, in which case it is applied to every integer
coefficient separately. The reductions are always done into the
symmetric interval $[-t/2,t/2)$.
$\log_a$ denotes the base-$a$ logarithm, and $\log$ always denotes
the base-2 logarithm. Table~1 below lists commonly used parameters,
and in some cases their corresponding names in SEAL 2.3.1.

\begin{table}[ht]
\centering
\renewcommand{\arraystretch}{1.3}
\begin{tabular}{|c|p{8.2cm}|p{3.5cm}|}
\hline
\textbf{Parameter} & \textbf{Description} &
\textbf{Name in SEAL (if applicable)} \\
\hline

$q$ &
Modulus in the ciphertext space (coefficient modulus) of the
form $q_1 \times \cdots \times q_k$, where $q_i$ are prime
&
\texttt{coeff\_modulus}
\\
\hline

$t$ &
Modulus in the plaintext space (plaintext modulus)
&
\texttt{plain\_modulus}
\\
\hline

$n$ &
A power of 2
&
\\
\hline

$x^n+1$ &
The polynomial modulus which specifies the ring $R$
&
\texttt{poly\_modulus}
\\
\hline

$R$ &
The ring $\mathbb{Z}[x]/(x^n+1)$
&
\\
\hline

$R_a$ &
The ring $\mathbb{Z}_a[x]/(x^n+1)$, i.e. same as the ring $R$
but with coefficients reduced modulo $a$
&
\\
\hline

$w$ &
A base into which ciphertext elements are decomposed during
relinearization
&
\\
\hline

$\log w$ &
&
\texttt{decomposition\_bit\_count}
\\
\hline

$\ell$ &
There are $\ell+1=\lfloor\log_w q\rfloor+1$ elements in each
component of each evaluation key
&
\\
\hline

$\delta$ &
Expansion factor in the ring $R$ ($\delta \leq n$)
&
\\
\hline

$\Delta$ &
Quotient on division of $q$ by $t$, or $\lfloor q/t\rfloor$
&
\\
\hline

$r_t(q)$ &
Remainder on division of $q$ by $t$, i.e.
$q=\Delta t+r_t(q)$, where $0\leq r_t(q)<t$
&
\\
\hline

$\chi$ &
Error distribution (a truncated discrete Gaussian distribution)
&
\\
\hline

$\sigma$ &
Standard deviation of $\chi$
&
\texttt{noise\_standard\_deviation}
\\
\hline

$B$ &
Bound on the distribution $\chi$
&
\texttt{noise\_max\_deviation}
\\
\hline

\end{tabular}
\caption{Notation used throughout this document.}
\label{tab:notation}
\end{table}